% The retailer's archive: a check of the mechanics on a second archive. Source comments follow the convention of results.tex: a "% src:" line names the report and key behind the one line that follows it.
\section{A Check of the Mechanics on a Second Archive}\label{sec:asos-results}
% src: $RUN6/backtest-asos-20261103-6bb5a15c5ff0.json header.tests_in_ledger 78, header.command (moot backtest --source asos, the same command as the headline archive's with the source changed)
The pipeline runs unchanged on a second public archive, the product experiments of an online retailer (Section~\ref{sec:setting}). That archive holds 78 experiments, which is too few to decide any claim. This appendix therefore checks that the mechanics run on other data, and it supports no finding.

Table~\ref{tab:fmt-asos} summarises Table~\ref{tab:main-asos} (Appendix~\ref{app:tables}), which applies the same protocol to the archive of the retailer, where the metric serves as the cell (Section~\ref{sec:setting}).
Every cell and the pooled archive are undecidable at their own gates.
% src: $RUN6/paper-extras-asos-20261103-6bb5a15c5ff0.json gates_by_stratum.strata.ALL own_gate.n_200 175, n_400 275, verdict undecidable at this n (metric_1 to metric_4: the same verdict at every rung, own gates n_200 from 175 to 250)
The smallest of those gates asks for 175 tests at a skill of 0.20.
% src: $RUN6/backtest-asos-20261103-6bb5a15c5ff0.json header.tests_scored 68
The count of scored tests for the archive is 68. This result was unavoidable, and the archive offers no basis for judging whether the findings on the headline archive carry over to product experiments.
% src: src/mootcourt/config.py SUPPORT_MIN 30; $RUN6/backtest-asos-20261103-6bb5a15c5ff0.json tables.abstention[cell=ALL] support 126, n 337, rate 0.490 (165 abstained rows); tables.origins abstention_rate 1.0, 1.0, 0.75 at the first three origins and 0 at the fourth and fifth; $RUN6/replay-asos-20261103-6bb5a15c5ff0.json weeks[].abstained_by_trigger.support 45, 48 and 36 in the first three quarters, 0 after
The reason the abstention rate runs high is mostly the support trigger, which withholds a card whenever its cell contains fewer than 30 visible rows (Section~\ref{sec:method}). The trigger fires at each of the first three origins and at no later origin, and it accounts for most of the rows that abstained.
For calibration, the figures the archive gives for the prior are descriptive: the tests are too few for us to treat them as a finding. Before adjustment its skill interval sits below zero, and after adjustment it spans zero. The expected calibration error of the prior lies beyond the rejection line for the clustered null while that of the base rate does not, and this is the single calibration figure that sets the two apart here.
% src: $RUN6/backtest-asos-20261103-6bb5a15c5ff0.json tables.aggregate[cell=ALL,rung=P1] coverage_80 0.7834 (coverage_80_lo 0.7175, coverage_80_hi 0.8438), coverage_90 0.8605 (coverage_90_lo 0.8042, coverage_90_hi 0.9114), p_sig_mean 0.1784, event_rate 0.1395, ece_percentile 98.8; [cell=ALL,rung=P0] p_sig_mean 0.1708, ece_percentile 64.9; tables/fmt-asos.tex
The coverage of its 80 and 90 percent intervals falls slightly short of the nominal levels, but the bootstrap interval around each coverage figure includes the nominal level, so under-coverage is not shown. On average it forecasts a significant positive result at a rate above the rate that was realized, and the base rate does the same, so the gap is not peculiar to the prior (Table~\ref{tab:fmt-asos}).

\begin{table}[t]\centering\scriptsize
% table: tab:main-asos; tables/fmt-asos.tex
% src: $RUN6/backtest-asos-20261103-6bb5a15c5ff0.json tables.aggregate[cell=ALL,rung=P1] family_size 55 (the intervals of tables/main-asos.tex), design_effect 3.82; tables/fmt-asos.tex
\caption{The archive of the retailer, taken from Table~\ref{tab:main-asos} (Appendix~\ref{app:tables}). Each row carries one figure. The adjusted interval is for the 55 intervals in Table~\ref{tab:main-asos}. The percentile for calibration is computed against the clustered null, and the design effect measures the inflation in variance that comes from the contrasts of a single test shifting together, which the null reproduces.}
\label{tab:fmt-asos}
\resizebox{\columnwidth}{!}{% generated by paper/build_tables.py from the report keys named above each row (results_b)
\begin{tabular}{lr}
\toprule
figure & value \\
\midrule
% src: $RUN6/backtest-asos-20261103-6bb5a15c5ff0.json header.scored_rows
scored contrasts & 337 \\
% src: $RUN6/backtest-asos-20261103-6bb5a15c5ff0.json header.tests_scored
experiments scored & 68 \\
% src: $RUN6/backtest-asos-20261103-6bb5a15c5ff0.json header.origins_scored
nominal quarterly origins & seven \\
% src: $RUN6/paper-extras-asos-20261103-6bb5a15c5ff0.json gates_by_stratum.strata.ALL.dependence.base_rate; $RUN6/paper-extras-asos-20261103-6bb5a15c5ff0.json gates_by_stratum.strata.ALL.own_gate.n_200; $RUN6/paper-extras-asos-20261103-6bb5a15c5ff0.json gates_by_stratum.strata.ALL.own_gate.n_400
own gate at skill 0.20 and 0.10, tests (base rate 0.14) & 175 and 275 \\
% src: $RUN6/backtest-asos-20261103-6bb5a15c5ff0.json tables.aggregate[cell=ALL,rung=P1].abstention_rate
abstention rate & 0.490 \\
% src: $RUN6/backtest-asos-20261103-6bb5a15c5ff0.json tables.abstention[cell=ALL].support
support trigger, rows at the early origins & 126 \\
% src: $RUN6/backtest-asos-20261103-6bb5a15c5ff0.json tables.aggregate[cell=ALL,rung=P1].brier_skill; $RUN6/backtest-asos-20261103-6bb5a15c5ff0.json tables.aggregate[cell=ALL,rung=P1].brier_skill_lo; $RUN6/backtest-asos-20261103-6bb5a15c5ff0.json tables.aggregate[cell=ALL,rung=P1].brier_skill_hi
P1 Brier skill (90\%) & -0.1073 (-0.2366, -0.0020) \\
% src: $RUN6/backtest-asos-20261103-6bb5a15c5ff0.json tables.aggregate[cell=ALL,rung=P1].adjusted_lo; $RUN6/backtest-asos-20261103-6bb5a15c5ff0.json tables.aggregate[cell=ALL,rung=P1].adjusted_hi
P1 Brier skill, adjusted interval & (-0.3258, 0.1113) \\
% src: $RUN6/backtest-asos-20261103-6bb5a15c5ff0.json tables.aggregate[cell=ALL,rung=P1].ece
expected calibration error & 0.127 \\
% src: $RUN6/backtest-asos-20261103-6bb5a15c5ff0.json tables.aggregate[cell=ALL,rung=P1].ece_percentile
calibration percentile & 98.8 \\
% src: $RUN6/backtest-asos-20261103-6bb5a15c5ff0.json tables.aggregate[cell=ALL,rung=P1].design_effect
design effect & 3.8 \\
% src: $RUN6/backtest-asos-20261103-6bb5a15c5ff0.json tables.aggregate[cell=ALL,rung=P1].coverage_80; $RUN6/backtest-asos-20261103-6bb5a15c5ff0.json tables.aggregate[cell=ALL,rung=P1].coverage_90
coverage of the 80 and 90 percent intervals & 0.78 and 0.86 \\
% src: $RUN6/backtest-asos-20261103-6bb5a15c5ff0.json tables.aggregate[cell=ALL,rung=P1].p_sig_mean; $RUN6/backtest-asos-20261103-6bb5a15c5ff0.json tables.aggregate[cell=ALL,rung=P1].event_rate
P1 mean forecast of the event and realized rate & 0.18 and 0.14 \\
% src: $RUN6/backtest-asos-20261103-6bb5a15c5ff0.json tables.aggregate[cell=ALL,rung=GBM+cal].brier_skill; $RUN6/backtest-asos-20261103-6bb5a15c5ff0.json tables.aggregate[cell=ALL,rung=GBM+cal].brier_skill_lo; $RUN6/backtest-asos-20261103-6bb5a15c5ff0.json tables.aggregate[cell=ALL,rung=GBM+cal].brier_skill_hi
boosted model, recalibrated Brier skill (90\%) & -0.0090 (-0.0756, 0.0446) \\
% src: $RUN6/backtest-asos-20261103-6bb5a15c5ff0.json tables.aggregate[cell=ALL,rung=local_eb+cal].brier_skill; $RUN6/backtest-asos-20261103-6bb5a15c5ff0.json tables.aggregate[cell=ALL,rung=local_eb+cal].brier_skill_lo; $RUN6/backtest-asos-20261103-6bb5a15c5ff0.json tables.aggregate[cell=ALL,rung=local_eb+cal].brier_skill_hi
local comparison set, recalibrated Brier skill (90\%) & 0.0150 (-0.0351, 0.0691) \\
% src: $RUN6/backtest-asos-20261103-6bb5a15c5ff0.json tables.decision_value[origin=ALL,budget=0.25,policy=p_win,rung=GBM].diff_random; $RUN6/backtest-asos-20261103-6bb5a15c5ff0.json tables.decision_value[origin=ALL,budget=0.25,policy=p_win,rung=GBM].diff_random_lo; $RUN6/backtest-asos-20261103-6bb5a15c5ff0.json tables.decision_value[origin=ALL,budget=0.25,policy=p_win,rung=GBM].diff_random_hi
boosted model, lift minus random order, quarter of tests (90\%) & 0.0420 (-0.0243, 0.1352) \\
% src: $RUN6/backtest-asos-20261103-6bb5a15c5ff0.json tables.decision_value[origin=ALL,budget=0.25,policy=p_win,rung=local_eb].diff_random; $RUN6/backtest-asos-20261103-6bb5a15c5ff0.json tables.decision_value[origin=ALL,budget=0.25,policy=p_win,rung=local_eb].diff_random_lo; $RUN6/backtest-asos-20261103-6bb5a15c5ff0.json tables.decision_value[origin=ALL,budget=0.25,policy=p_win,rung=local_eb].diff_random_hi
local comparison set, the same (90\%) & 0.0044 (-0.0660, 0.0934) \\
% src: $RUN6/backtest-asos-20261103-6bb5a15c5ff0.json tables.decision_value[origin=ALL,budget=0.25,policy=p_win,rung=GBM].diff_matched; $RUN6/backtest-asos-20261103-6bb5a15c5ff0.json tables.decision_value[origin=ALL,budget=0.25,policy=p_win,rung=local_eb].diff_matched
both, lift minus the matched order, quarter of tests & 0.000 and 0.000 \\
% src: $RUN6/paper-extras-asos-20261103-6bb5a15c5ff0.json impressions_effect_correlation.corr_log_impressions_abs_z
exposure stand-in, correlation with $|$standardized effect$|$ & 0.215 \\
% src: $RUN6/paper-extras-upworthy-20261106-2dc461f52ea4.json impressions_effect_correlation.corr_log_impressions_abs_z
the same correlation on the headline archive & 0.087 \\
\bottomrule
\end{tabular}
}
\end{table}

% src: $RUN6/backtest-asos-20261103-6bb5a15c5ff0.json tables.p2_selection_by_cell share_lambda_zero 1.00 in all four cells; header.challengers GBM, GBM+cal, local_eb, local_eb+cal (no text classifier); tables.aggregate[cell=ALL] brier_skill GBM -0.7184, local_eb -0.3832, GBM+cal -0.0090 (-0.0756, 0.0446), local_eb+cal 0.0150 (-0.0351, 0.0691), GBM and local_eb brier_skill_hi below zero and adjusted_lo, adjusted_hi spanning zero; tables.decision_value origin=ALL budget=0.25 policy=p_win diff_random GBM 0.0420 (-0.0243, 0.1352), local_eb 0.0044 (-0.0660, 0.0934), diff_matched 0.0 (src/mootcourt/decision_value.py _matched_value)
On every row P2 comes out equal to P1: the admission guard of Section~\ref{sec:method} admits no precedent weight in any cell at any origin. The challengers that do run here are two: the boosted model and the local comparison set. The text classifier requires headlines, and this archive has none. In raw form the point estimates of both challengers fall far below the base rate, their intervals lie below zero before adjustment and span zero after it, and once recalibrated both intervals span zero. As with the prior, we treat none of this as a finding. Against random order their captured lift spans zero, and it equals zero exactly against the matched order, which on this archive gives back the value of the chosen set itself (Appendix~\ref{app:asos}).
